\section{Future Work}
\label{sec:future}

In this section, we discuss several unresolved research challenges associated with honesty and provide insights into potential research avenues.

\paragraph{Objective or Subjective.} A central debate in current research on the honesty of LLMs revolves around whether honesty should be considered a subjective or objective concept. Askell et al. (2021); Kadavath et al. (2022) describe honesty as the ability to provide accurate information along with calibrated confidence reflecting the correctness of its answers. In contrast, Evans et al. (2021); Lin et al. (2022a) view honesty as the model's ability to express its own beliefs. The former takes an objective approach, aligning with world knowledge, while the latter adopts a subjective perspective, focusing on the model's internal state. The objective perspective better suits human needs, as people generally value accurate and truthful information. However, this approach presents challenges for optimizing models because there is often a gap between the model's knowledge and world knowledge, necessitating additional supervision during training to \emph{distinguish between true and false information}. This becomes particularly challenging for future superhuman models, where human oversight may be limited to providing only weak supervision (Burns et al., 2023a). Conversely, the subjective perspective might require less supervision, as it primarily focuses on the model's ability to express its own knowledge, and the challenge lies in \emph{differentiating between what they know and unknown knowledge}. Nonetheless, even if a model can fully articulate its knowledge, problems arise when that knowledge is incorrect. Both the objective and subjective perspectives have distinct advantages and challenges, and resolving this debate is crucial for further research progress.

\paragraph{Knowledge Identification.} As stated in previous sections, knowledge identification has been the primary challenge in both evaluation and methodological approaches. However, the exact definitions of what should be considered known or unknown remain unclear. Existing research typically follows two mainstream strategies: a supervised approach, which distinguishes them based on the correctness of responses, and an unsupervised approach, which differentiates them based on the uncertainty in responses. Both strategies depend on the external expression of LLMs, but what if these models struggle to express what they know? Indeed, several studies have highlighted the discrepancy between LLMs' internal knowledge and what they express (Liu et al., 2023; Li et al., 2024a). Ignoring this issue could limit the potential of LLMs to fully express their knowledge. Therefore, in addition to focusing on external expression, future research could explore techniques that utilize the inherent knowledge LLMs possess to differentiate between known and unknown information (Burns et al., 2023b; Hernandez et al., 2023; Wang et al., 2024b). For instance, researchers could aim to recover the knowledge embedded in model parameters or present in the contexts even in cases when LLMs fail to express it faithfully in their outputs (Burns et al., 2023b).

\paragraph{Honesty in Instruction-following.} Current research on honesty focuses primarily on knowledge-intensive question answering, particularly those involving short-form answers, while largely overlooking instruction-following scenarios that are more desired in real-world applications. Instruction-following differs from question answering as it requires broader capabilities and typically involves long-form generation. To address this gap, future research can establish evaluation methods and benchmarks to assess the honesty of LLMs in instruction-following. Additionally, they can explore methods for improvement, such as prompting, supervised fine-tuning, and reinforcement learning.

\paragraph{Honesty on In-context Knowledge.} As noted in \S\ref{sec:what_honesty}, LLMs possess two types of knowledge: internal parametric knowledge acquired through training and external in-context knowledge. While most existing research emphasizes the honesty of parametric knowledge, the honesty of in-context knowledge has received less attention. However, in real-world applications of LLMs, in-context knowledge also plays a vital role in generation, particularly in retrieval-augmented and long-context scenarios. Therefore, we advocate for future research to devote more attention to the honesty of in-context knowledge.

\paragraph{Honesty in Various Models.} Research on honesty has primarily focused on transformer decoder-based LLMs. However, other popular models also deserve attention, including multimodal LLMs such as GPT-4V (Achiam et al., 2023) and Gemini (Team et al., 2023), models with novel architectures such as Mamba (Gu \& Dao, 2023), and compressed models using compression techniques such as quantization and pruning (Wan et al., 2024b). We believe that these models also merit further exploration in future research.
